\section{Mathematical background}


\subsection{The forward operator}
\input{figures/forward_operator.tex}
Let $\Omega \subset \mathbb R^2$ be a bounded domain with boundary $\partial \Omega$, and consider the boundary value problem (BVP) in Equation~\eqref{eq:pde}. The coefficient $\sigma = \sigma(x) > 0$ denotes the diffusion coefficient and determines the rate and direction of diffusion in the problem. It may be either scalar-valued or matrix-valued. In the latter case, we assume that $\sigma(x)$ is symmetric and positive definite (see Appendix~\ref{sec:positive_definiteness}) to keep the problem elliptic. The reaction coefficient $c = c(x) > 0$ describes the absorption or decay within the problem. Finally, $f = f(x)$ denotes the source term, which is the unknown input of the inverse problem.
\nomenclature[A]{$\Omega$}{Computational domain}
\nomenclature[A]{$\partial\Omega$}{Boundary of the domain}
\nomenclature[B]{$\sigma,\, \sigma(x)$}{Diffusion coefficient}
\nomenclature[B]{$c,\, c(x)$}{Reaction coefficient}

Equation~\eqref{eq:pde} is a linear diffusion-reaction problem with Neumann boundary conditions, and defines a mapping
\begin{equation}
    \mathcal S: L^2(\Omega) \rightarrow H^1(\Omega), \quad \mathcal{S} f = u.
\end{equation}
Here, $L^2(\Omega)$ is the space of square integrable functions and $H^1(\Omega) \subset L^2(\Omega)$ the first-order Sobolev space, see Appendix~\ref{appendix:function_spaces}
\nomenclature[C]{$\mathcal{S}$}{PDE operator}
\nomenclature[D]{$L^2(\Omega)$}{Space of square integrable functions}
\nomenclature[D]{$H^1(\Omega)$}{First-order Sobolev space}

To model boundary observations, we introduce the trace operator
\begin{equation}\label{eq:trace}
    \mathcal T: H^1(\Omega) \rightarrow L^2(\partial\Omega), \quad \mathcal{T} u = u|_{\partial\Omega}
\end{equation}
\nomenclature[C]{$\mathcal{T}$}{Trace operator}
The term $\mathcal T u$ represents the observed values of $u$ on the boundary of $\Omega$.

The forward operator is then defined as the composition
\begin{equation}\label{eq:K}
    \mathcal K = \mathcal T \circ \mathcal S, \quad \mathcal{K} f = u|_{\partial\Omega},
\end{equation}
\nomenclature[C]{$\mathcal{K}$}{Forward operator}
which is a mapping from the source function $f$ to the observed output $u|_{\partial\Omega} = \mathcal K f$ on the boundary $\partial \Omega$. See Figure~\ref{fig:forward_operator}.

The associated inverse problem can now be formulated as follows: given boundary data $u|_{\partial\Omega}$, determine a source function $f$ such that
\begin{equation}
    \mathcal{K} f = u|_{\partial\Omega}.
\end{equation}
In order to solve such a problem computationally, the problem must be discretized. In this thesis, the discretization is done using the finite element method (FEM), described in the following sections.


\subsection{The finite element method}
\input{figures/triangulation.tex}
Consider the PDE in Eq.~\eqref{eq:pde}. Multiplying the equation by a test function $v \in H^1(\Omega)$, integrating over $\Omega$, and applying Green's identity yields the weak formulation: find $u \in H^1(\Omega)$ such that
\begin{equation}\label{eq:variational_formulation}
    \int_{\Omega} (\sigma \nabla u \cdot \nabla v + k u v) \, dx = \int_{\Omega} f v \space \, dx,
    \text{ for all } v \in H^1(\Omega).
\end{equation}
For details on the derivation, see Section~4 in Larson and Bengzon \cite{larson_finite_2013}.

Let $\mathcal M$ be a triangulation of $\Omega$ made up of $N$ nodes, with $N_b$ boundary nodes, see Fig.~\ref{fig:triangulation} as an example. Introduce the finite element space $V_h \subset H^1(\Omega)$,
\begin{equation}\label{eq:fem_space}
    V_h = \{v \in C^0(\Omega) : v|_M \text{ is linear for all } M \in \mathcal M \}.
\end{equation}
Let $\{\phi_1, \dots, \phi_N\}$ denote a basis for $V_h$, so that
\begin{equation}
    V_h = \operatorname{span}\{ \phi_1, \dots, \phi_N \}.
\end{equation}
The basis functions $\phi_i$ are piecewise linear hat functions associated with the nodes of the mesh (see \cite[p.~51]{larson_finite_2013}).

\nomenclature[A]{$\mathcal{M}$}{Triangulation (mesh)}
\nomenclature[A]{$N$}{Number of mesh nodes}
\nomenclature[A]{$N_b$}{Number of boundary mesh nodes}
\nomenclature[D]{$V_h$}{Finite element space}

The finite element approximation of \eqref{eq:variational_formulation} then reads: 
find $u_h \in V_h$ such that
 \begin{equation}\label{eq:fem_formulation}
    \int_{\Omega} (\sigma \nabla u_h \cdot \nabla v + k u_h v) \, dx = \int_{\Omega} f v \space \, dx,
    \text{ for all } v \in V_h.
\end{equation}


\subsection{The discrete forward operator}
To obtain a discrete representation of the forward operator 
$\mathcal{K} = \mathcal{T} \circ \mathcal{S}$, 
we discretize both the PDE operator $\mathcal{S}$ and the trace operator $\mathcal{T}$ using the finite element space $V_h$.

Let $f_h, u_h \in V_h$ denote the finite element approximations of $f$ and $u|_{\partial\Omega}$, respectively. Since $V_h = \operatorname{span}\{\phi_1,\dots,\phi_N\}$, they can be written as
\begin{equation}\label{eq:uh_and_fh}
    f_h = \sum_{j=1}^{N} x_j \phi_j, \quad
    u_h = \sum_{j=1}^{N} u_j \phi_j, 
\end{equation}
with coefficient vectors $\v u = (u_1,\dots,u_N)^T$ 
and $\v x = (x_1,\dots,x_N)^T$.

Inserting \eqref{eq:uh_and_fh} into the finite element formulation \eqref{eq:fem_formulation} and testing with the basis functions 
$\{\phi_i\}_{i=1}^N$ yields the linear system
\begin{equation}\label{eq:matrix_system}
    A \mathbf{u} = M \mathbf{x},
\end{equation}
where $A \in \mathbb{R}^{N \times N}$ is the stiffness matrix and $M \in \mathbb{R}^{N \times N}$ the mass matrix, see Appendix~\ref{appendix:stiffness}.
\nomenclature[E]{$M$}{Volume mass matrix}%
Equation~\eqref{eq:matrix_system} implies the discrete PDE operator $S$ takes the form
\begin{equation}
    S = A^{-1} M \in \mathbb{R}^{N \times N}
\end{equation}
Note that $S$ is a mapping of coefficients vectors: mapping $\v x$ to $\v u$.

The trace operator is discretized by a matrix 
$T \in \mathbb{R}^{N_b \times N}$ that extracts the boundary coefficients. If $\mathbf{y}$ denotes the vector of boundary coefficients, then
\begin{equation}
    \v y = T \v u.
\end{equation}

Combining the discrete PDE and trace operators yields the matrix representation of the forward operator,
\begin{equation}\label{eq:discrete_K}
    K = T S = T A^{-1} M \in \mathbb{R}^{N_b \times N}.
\end{equation}
\nomenclature[E]{$K$}{Discrete forward operator}%


\subsection{Discrete inner product and norm}
Let $V_h$ be the finite element space introduced in Eq.~\eqref{eq:fem_space} and let $f_h, g_h \in V_h$, with coefficient vectors $\v x, \v y \in \mathbb{R}^N$, respectively. The $L^2(\Omega)$ inner product (see Eq.~\eqref{eq:L2_inner_product}) of $f_h$ and $g_h$, can be written as
\begin{equation}
    \langle f_h, g_h \rangle_{L^2(\Omega)} = \v x^T M \v y,
\end{equation}
where $M$ is the mass matrix. 

This induces a discrete inner product on $\mathbb{R}^N$, with associated norm, defined by
\begin{equation}
    \langle \v x, \v y \rangle_M = \v x^T M \v y,
    \quad
    \| \v x \|_M = \sqrt{\v x^T M \v x}.
\end{equation}

Similarly, let $\v u, \v v \in \mathbb{R}^{N_b}$ 
denote coefficient vectors corresponding to the traces 
$u_h|_{\partial\Omega}$ and $v_h|_{\partial\Omega}$. The $L^2(\partial\Omega)$ inner product induces a discrete inner product and norm on $\mathbb{R}^{N_b}$, given by
\begin{equation}
    \langle \mathbf{u}, \mathbf{v} \rangle_{M_{\partial}}
    = \mathbf{u}^T M_{\partial} \mathbf{v},
    \quad
    \| \v  u \|_{M_{\partial}}
    =
    \sqrt{ \v  u^T M_{\partial} \v u},
\end{equation}
where $M_\partial$ denotes the boundary mass matrix, see Appendix~\ref{appendix:stiffness}.
\nomenclature[E]{$M_\partial$}{Boundary mass matrix}%

%In practice, this system is often **ill-posed** due to the large null space of $K$ and the accumulation of numerical errors in $A^{-1}$. As a result, naive solutions of Eq.~\eqref{eq:inverse_problem} are highly sensitive to noise in $\v y$, motivating the need for regularization techniques. In the following section, we discuss classical Tikhonov regularization and introduce the weighted approach that will be explored in this thesis.


\subsection{Tikhonov regularization}
% Introduce general min problem, and normal equation.
% Then introduce Tikhonov 
% Then introduce general Tikhonov
Let $K$ be an $N_b \times N$ matrix, and consider the problem
\begin{equation}\label{eq:tikhonov}
    \min_{x \in \mathbb R^N} \| K x - y \|^2 + \lambda^2 \| W x \|^2,
\end{equation}
where $\lambda \ge 0$ is the regularization parameter and $W \in \mathbb{R}^{N \times N}$ are the regularization weights. The solution $x_\lambda$ of Eq.~\eqref{eq:tikhonov} is known as the Tikhonov solution of $K x = y$. The term $\| W x \|^2$ punishes the weighted norm $x$ preventing the norm of $x$ from exploding, which often is a problem in ill-posed problems \textcolor{red}{(reference)}.


\subsection{Randomized SVD}
The rSVD is a technique used to compute a rank-$k$ approximation of an $m\times n$ matrix without computing the SVD. The core idea is to use random sampling to find a low-rank matrix $Q$ which approximates the range of $A$, project $A$ onto $\operatorname{span}(Q)$, and compute the SVD of the projection. The projection $B=Q^TA$ will, with high probability, capture the dominant actions of $A$, and the SVD of $B$ will yield a good approximation \cite{halko_finding_2011}. This procedure is summarized in Algorithm \ref{alg:rsvd}. The construction of an orthonormal basis $Q$ of $Y$ is done through a QR factorization. Note that the approximation $A\approx QQ^T A$ and the approximate factorization $A\approx U\Sigma V^T$ produced by rSVD have the same error. This can be seen by observing that $QQ^TA = QB = Q\tilde U\Sigma V^T$, and setting $U=Q\tilde U$.
\begin{algorithm}[!htbp]
\caption{rSVD}
\label{alg:rsvd}
\begin{algorithmic}[1]
    \Require $A\in\mathbb R^{m\times n}$, target rank $k$, oversampling parameter $p$.
    \Ensure Approximate rank-$k$ factorization $U\Sigma V^T$.
    \State Generate a random $n\times(k+p)$ test matrix $\Omega$.
    \State Compute $Y = A\Omega$.
    \State Form an orthonormal basis $Q$ of $Y$.
    \State Set $B = Q^TA$.
    \State Compute the SVD $B = \tilde U\Sigma V^T$.
    \State Set $U = Q\tilde U$. 
\end{algorithmic}
\end{algorithm}